% Put the project assessment here
\section{Project Assessment}

This chapter outlines the evaluation of the VERACISM once it will be developed. The aim of the evaluation process is to assure that the system is operational, secure, usable, and deployable in the International School Manila (ISM) environment. In the proposal phase the emphasis is on the proposed approach to testing; the tests to be used; and the metrics to be gathered. These actual tests will produce test results only after implementation.

\subsection{Functional Testing}

The purpose of this testing was to ensure that the features that were implemented under VERACISM were functioning as they were designed to do so based on the system requirements. These testing procedures included validating uploaded files, checking for malware, hashing uploaded content, generating CIDs, downloading content with control, auditing and dashboard views, tenant-scoped access and public verification.

A \textbf{hybrid approach} of
\textbf{automated black-box API testing} and \textbf{manual role-based
functional testing}. Automated execution: Automated using Newman in Postman collections with the staging environment. These automated tests: tested the v1 public API workflow, some end-to-end user journeys, including IT Office upload, Registrar upload, student and external verification, protected download behaviour; and Cross-role upload to verify flow.

For the sake of manual functional testing, a few modules and flows that are exposed on the user interface needed manual testing, some Bearer JWT protected internal endpoints are involved, a few tests using QR code with a mobile device, a few reports on the dashboard, a couple of access checks followed by tenant scope and some gates managed by the administration program that needs role interaction.

Manual functional procedures were developed and followed for the System Administrator, IT Office / tenant user, Administration Office / Registrar, Student, External Auditor / Educator and for the end-to-end scenarios. These procedures included tenant approval, user management, API key updates, configuration changes, audit log review/export, health monitoring, QR display and scan verification, dashboard status checking, report deletion or immutability checks, comprehension checks, access boundary checks across tenants, and multi-network verification checks.

The primary criteria used in functional testing included correct HTTP status
codes, correct response structures, presence and format of expected fields,
correct role and tenant access enforcement, correct UI behaviour where
applicable, successful completion of required user journeys, and retention of
the required audit evidence. Supporting automated and manual functional
testing records will be retained as appendix evidence.

\subsubsection{Functional test success criteria}

To make the functional assessment explicit and reproducible, each test area is
defined with measurable criteria for a \textit{Pass} result. These criteria
focus on correct outputs, correct security behaviour, correct logging,
observable handling of failures, and coverage of both automated and manual
functional modules.

% ---------------------------
% Functional table (Part 1)
% ---------------------------
\begin{table}[H]
  \centering
  \smalltable
  \footnotesize
  {%
  \setlength{\tabcolsep}{5pt}%
  \renewcommand{\arraystretch}{1.35}%

  \begin{tabular}{|>{\raggedright\arraybackslash}p{0.16\textwidth}|
                  >{\raggedright\arraybackslash}p{0.16\textwidth}|
                  >{\raggedright\arraybackslash}p{0.24\textwidth}|
                  >{\raggedright\arraybackslash}p{0.12\textwidth}|
                  >{\raggedright\arraybackslash}p{0.16\textwidth}|
                  >{\raggedright\arraybackslash}p{0.08\textwidth}|}
    \hline
    \textbf{Planned TC ID(s)} &
    \textbf{Method} &
    \textbf{Success criteria for pass} &
    \textbf{Target} &
    \textbf{Evidence} &
    \textbf{Result} \\
    \hline\hline

    TC-ITO-001, TC-REG-001 &
    Automated &
    Valid authenticated uploads are accepted, stored in temporary storage,
    initial metadata is recorded, and a tracking identifier is returned. &
    Accepted upload behaves as specified &
    API logs, tracking ID, temp storage record &
    Pass / Fail \\
    \hline

    TC-ITO-002, TC-REG-002 &
    Automated &
    Invalid, oversized, unsupported, or unauthorized uploads are rejected and
    do not create stored artifacts. &
    Rejected cases behave as specified &
    Error response, logs, absence of stored artifact &
    Pass / Fail \\
    \hline

    TC-ITO-003, TC-REG-003 &
    Automated &
    Retrieved report metadata, student association, SHA-256 hash, and CID are
    consistent with the uploaded submission. &
    Stored values match submitted values &
    API response, metadata comparison, lookup result &
    Pass / Fail \\
    \hline

    TC-ITO-004, TC-REG-004, TC-STU-001, TC-EXT-001 &
    Automated + Manual &
    Successful uploads produce a usable CID and verification path, and QR-based
    verification resolves correctly through the public portal. &
    CID present and QR resolves correctly &
    API response, UI screenshot, mobile scan proof &
    Pass / Fail \\
    \hline

    TC-REG-005, TC-ADM-004 &
    Manual &
    Dashboard and audit views show the correct upload, anchoring, and audit
    information, and exports work in the required formats. &
    Statuses visible and export succeeds &
    UI screenshots, export file, audit review notes &
    Pass / Fail \\
    \hline

  \end{tabular}
  }%
  \captionsetup{type=table}
  \caption{Functional test success criteria with automated and manual coverage (Part 1)}
  \label{tab:functional-success-criteria-a}
\end{table}

% ---------------------------
% Functional table (Part 2)
% ---------------------------
\begin{table}[H]
  \centering
  \smalltable
  \footnotesize
  {%
  \setlength{\tabcolsep}{4pt}%
  \renewcommand{\arraystretch}{1.35}%

  \begin{tabular}{|>{\raggedright\arraybackslash}p{0.17\textwidth}|
                  >{\raggedright\arraybackslash}p{0.14\textwidth}|
                  >{\raggedright\arraybackslash}p{0.29\textwidth}|
                  >{\raggedright\arraybackslash}p{0.13\textwidth}|
                  >{\raggedright\arraybackslash}p{0.15\textwidth}|
                  >{\raggedright\arraybackslash}p{0.08\textwidth}|}
    \hline
    \textbf{Planned TC ID(s)} &
    \textbf{Method} &
    \textbf{Success criteria for pass} &
    \textbf{Target} &
    \textbf{Evidence} &
    \textbf{Result} \\
    \hline\hline

    TC-ADM-001, TC-ADM-002, TC-ADM-003 &
    Manual &
    Tenant approval, user and role management, API key management, and
    critical configuration changes can be completed through the admin module
    and reflect correctly in the system. &
    Admin actions succeed and are reflected correctly &
    UI screenshots, admin logs, review notes &
    Pass / Fail \\
    \hline

    TC-ADM-005, TC-ITO-006, TC-STU-005, TC-EXT-004 &
    Manual + Automated &
    Restricted views and actions are accessible only to authorized roles.
    Unauthorized attempts are rejected, and tenant boundaries are enforced. &
    0 unauthorized privileged actions succeed &
    Access logs, route checks, API rejection results &
    Pass / Fail \\
    \hline

    TC-REG-006 &
    Manual &
    Finalized report content cannot be casually edited through normal user
    paths, and any delete behaviour follows the intended soft-delete or
    policy-controlled workflow. &
    No improper content modification path exposed &
    UI inspection, code review note, behaviour check &
    Pass / Fail \\
    \hline

  \end{tabular}
  }%
  \captionsetup{type=table}
  \caption{Functional test success criteria with automated and manual coverage (Part 2)}
  \label{tab:functional-success-criteria-b}
\end{table}

% ---------------------------
% Functional table (Part 3)
% ---------------------------
\begin{table}[H]
  \centering
  \smalltable
  \footnotesize
  {%
  \setlength{\tabcolsep}{4pt}%
  \renewcommand{\arraystretch}{1.35}%

  \begin{tabular}{|>{\raggedright\arraybackslash}p{0.17\textwidth}|
                  >{\raggedright\arraybackslash}p{0.14\textwidth}|
                  >{\raggedright\arraybackslash}p{0.29\textwidth}|
                  >{\raggedright\arraybackslash}p{0.13\textwidth}|
                  >{\raggedright\arraybackslash}p{0.15\textwidth}|
                  >{\raggedright\arraybackslash}p{0.08\textwidth}|}
    \hline
    \textbf{Planned TC ID(s)} &
    \textbf{Method} &
    \textbf{Success criteria for pass} &
    \textbf{Target} &
    \textbf{Evidence} &
    \textbf{Result} \\
    \hline\hline

    TC-STU-002, TC-STU-003, TC-EXT-002, TC-E2E-001, TC-E2E-002 &
    Automated &
    Valid CID verification shows the correct result and invalid or fabricated
    CID values return a clear warning without internal error leakage. &
    Correct result and message shown for all cases &
    Verification response, screenshots, timing capture &
    Pass / Fail \\
    \hline

    TC-STU-004, TC-EXT-003, TC-E2E-003 &
    Manual / Hybrid &
    Non-technical users can understand the verification result, repeated checks
    from different networks remain consistent, and unavailable content cases
    show clear warnings without exposing sensitive details. &
    Result understandable and behaviour consistent &
    Observation notes, network comparison notes, warning screen proof &
    Pass / Fail \\
    \hline

  \end{tabular}
  }%
  \captionsetup{type=table}
  \caption{Functional test success criteria with automated and manual coverage (Part 3)}
  \label{tab:functional-success-criteria-c}
\end{table}

\subsubsection{Functional test result classification}

For functional testing, each planned test case will be classified using one of
the following outcome labels.

\begin{table}[H]
  \centering
  \smalltable
  \footnotesize
  {%
  \setlength{\tabcolsep}{6pt}
  \renewcommand{\arraystretch}{1.3}
  \begin{tabular}{|>{\raggedright\arraybackslash}p{0.22\textwidth}|
                  >{\raggedright\arraybackslash}p{0.68\textwidth}|}
    \hline
    \textbf{Outcome label} & \textbf{Definition} \\
    \hline\hline

    \textbf{Pass} &
    All expected behaviours, outputs, access controls, and required evidence
    are observed, and the result matches the stated success criteria. \\
    \hline

    \textbf{Fail} &
    One or more expected behaviours, outputs, controls, or validations are not
    met, or incorrect behaviour is observed during execution. \\
    \hline

    \textbf{Partial} &
    The feature is only partly implemented, the flow is only partly testable,
    or the observed behaviour satisfies only part of the expected result. \\
    \hline

    \textbf{Not Implemented} &
    The required feature, control, or supporting logic is not yet available in
    the current build, so the planned test case cannot be fully executed. \\
    \hline

  \end{tabular}
  }%
  \captionsetup{type=table}
  \caption{Functional test result classification}
  \label{tab:functional-result-classification}
\end{table}

Accordingly, the functional testing tables define the conditions required for a
\textit{Pass}, while actual execution results are recorded separately using the
classification labels in Table~\ref{tab:functional-result-classification}.

\subsection{Integration Testing}

Integration testing focused on verifying that the major VERACISM components
worked correctly when combined into a complete operational pipeline. The
integration scope covered the interaction between the .NET 8 backend API,
SQL Server, ClamAV, YARA, Lighthouse/IPFS storage, and the public
verification process.

The integration assessment was performed through automated API-based tests in
the staging environment. These tests targeted the public API v1 endpoints
through API key authentication and covered the main integrated operations of
the platform, including upload and scan processing, report retrieval by CID,
public verification, controlled file download, and storage-status checking.
Both valid and invalid cases were included to assess API contract behaviour,
error handling, response consistency, and subsystem coordination.

The integration design also included a test for valid uploads to go to the scan, hash, metadata and storage points and invalid or unauthorized requests being blocked and not creating unnecessary records. More focus was put on tenant isolation, consistency of uploaded metadata retrieved metadata, and public versus protected endpoint behaviour.

The primary criteria used in integration testing included end-to-end request
completion, correctness of inter-service outputs, validation of report and CID
data across steps, correct rejection of invalid requests, and preservation of
system boundaries across tenants and access levels.

\subsection{Performance Testing}

The customer was interested in testing of the VERACISM responsiveness and stability against the expected verification traffic in the Staging environment. The public CID verification endpoint was focused for this as this is a key aspect of how the platform is used by students, external verifiers and partner institutions needing to make a rapid integrity check.

The performance run was performed on Postman Performance Testing and connected to Postman API staging environment, which is the VERACISM API. A total of 100 virtual users were used for 5 minutes in the test profile. The test scenario was aimed at the genuine public verification request with the capture of the request volume, throughput, response time, percentile response times and captured transport level errors during the run.

This assessment criteria covered the following: (1) Minimum error rate, (2) Stable throughput under concurrent traffic, (3) Average response time (within acceptable limits), and (4) Endpoint availability (sustained throughout the test window) during the test session. The metrics captured are stored as performance evidence and will be compared to and discussed with in relation to other captured metrics in the results chapter.

\subsection{User Acceptance Testing}

User acceptance testing will focus on the primary users of VERACISM: tenant administrators, system administrators, IT integrators and external verifiers through the public portal. The test scenarios will come from the essential functional flows and capabilities of the system, such as dashboard review, report and file management, tenant administration, setting the security configuration, reviewing the audit-log, student and integrator administration, public verification, and API-based integration workflows.

User Acceptance Assessment will be conducted with task-based scripts based on system Acceptance criterion. Each UAT case will outline the role of the participant, any related planned test case(s), the task the participant was required to perform, and the outcome and expected result relevant to that task. The participants will carry out the tasks set, and observers will document whether or not the tasks were completed, how long they took, any defects found and the task blocker observed during the carrying out of the tasks. Results will be
recorded using objective completion outcomes such as \textit{Passed},
\textit{Partially Completed}, or \textit{Failed}, supported by notes and
evidence captured during the session.

\begin{table}[htbp]
  \centering
  \smalltable
  \footnotesize
  {%
  \setlength{\tabcolsep}{8pt}
  \renewcommand{\arraystretch}{1.3}
  \begin{tabular}{|>{\raggedright\arraybackslash}p{3.3cm}|
                  >{\centering\arraybackslash}p{2.0cm}|
                  >{\raggedright\arraybackslash}p{6.6cm}|}
    \hline
    \textbf{Planned user group} &
    \textbf{Target \#} &
    \textbf{Representative test tasks} \\
    \hline\hline

    System Administrator (Super Admin) &
    1--2 &
    Review dashboard analytics and service health; manage reports and tenant
    records; configure YARA and ClamAV settings; manage API keys; review and
    export audit logs. \\
    \hline

    Tenant Administrator (Administration Office / ISM offices) &
    2-3 &
    Review dashboard summaries; upload and manage files; review version
    history; copy CID and download files; manage students; manage integrator
    users and API keys. \\
    \hline

    IT Integrators / Developers &
    2-3 &
    Submit reports through the API; retrieve report details; download report
    files; check file processing or status results; handle invalid or rejected
    API requests correctly. \\
    \hline

    External verifiers (students, auditors, public user) &
    4--5 &
    Verify records through QR code or CID entry using the public portal and
    confirm that valid, invalid, tampered, or unavailable cases display the
    correct result. \\
    \hline
  \end{tabular}
  }%
  \captionsetup{type=table}
  \caption{Planned user-testing participants and task coverage}
  \label{table:user-testing-plan}
\end{table}

\subsubsection{User acceptance success criteria}

To make the user acceptance assessment explicit and reproducible, each UAT
case is defined with measurable criteria for a \textit{Passed} result. These
criteria focus on successful task completion, achievement of the expected
result, absence of blocking defects, and the presence of sufficient
supporting evidence for the recorded outcome.

% =========================================================
% System Administrator UAT
% =========================================================
\begin{table}[H]
  \centering
  \smalltable
  \scriptsize
  {%
  \setlength{\tabcolsep}{3pt}
  \renewcommand{\arraystretch}{1.18}
  \begin{tabular}{|>{\raggedright\arraybackslash}p{0.09\textwidth}|
                  >{\raggedright\arraybackslash}p{0.15\textwidth}|
                  >{\raggedright\arraybackslash}p{0.22\textwidth}|
                  >{\raggedright\arraybackslash}p{0.30\textwidth}|
                  >{\raggedright\arraybackslash}p{0.14\textwidth}|
                  >{\centering\arraybackslash}p{0.10\textwidth}|}
    \hline
    \textbf{UAT ID} &
    \textbf{Linked TC ID(s)} &
    \textbf{Task} &
    \textbf{Success criteria for passed result} &
    \textbf{Evidence} &
    \textbf{Result} \\
    \hline\hline

    UAT-01 &
    TC-ADM-001, TC-ADM-002, TC-ADM-004 &
    Review administrator dashboard and core records. &
    Dashboard data, recent records, and report information are visible and can
    be reviewed without access or display errors. &
    Screenshots, viewed records &
    Passed / Partially Completed / Failed \\
    \hline

    UAT-02 &
    TC-ADM-001, TC-ADM-002, TC-ADM-003 &
    Manage tenants, users, and API keys. &
    Tenant records, user assignments, and API key actions are completed
    successfully and reflected correctly in the interface and audit trail. &
    Screenshots, audit entries &
    Passed / Partially Completed / Failed \\
    \hline

    UAT-03 &
    TC-ADM-003, TC-ADM-004, TC-ADM-005 &
    Manage security and monitoring settings. &
    YARA and ClamAV related actions complete successfully, system health is
    visible, and audit logs can be reviewed or exported as intended. &
    Screenshots, logs, export output &
    Passed / Partially Completed / Failed \\
    \hline

  \end{tabular}
  }%
  \captionsetup{type=table}
  \caption{User acceptance success criteria for System Administrator}
  \label{tab:uat-success-admin}
\end{table}

% =========================================================
% Tenant Administrator UAT
% =========================================================
\begin{table}[H]
  \centering
  \smalltable
  \scriptsize
  {%
  \setlength{\tabcolsep}{3pt}
  \renewcommand{\arraystretch}{1.18}
  \begin{tabular}{|>{\raggedright\arraybackslash}p{0.09\textwidth}|
                  >{\raggedright\arraybackslash}p{0.15\textwidth}|
                  >{\raggedright\arraybackslash}p{0.22\textwidth}|
                  >{\raggedright\arraybackslash}p{0.30\textwidth}|
                  >{\raggedright\arraybackslash}p{0.14\textwidth}|
                  >{\centering\arraybackslash}p{0.10\textwidth}|}
    \hline
    \textbf{UAT ID} &
    \textbf{Linked TC ID(s)} &
    \textbf{Task} &
    \textbf{Success criteria for passed result} &
    \textbf{Evidence} &
    \textbf{Result} \\
    \hline\hline

    UAT-04 &
    TC-REG-001, TC-REG-003, TC-REG-004, TC-REG-005 &
    Upload and manage report files. &
    Valid file operations complete successfully, reports are associated
    correctly, CID and QR data are available, and file records appear with the
    correct status. &
    Upload output, screenshots, resulting record &
    Passed / Partially Completed / Failed \\
    \hline

    UAT-05 &
    TC-REG-002, TC-REG-006 &
    Handle invalid file actions and restricted modification paths. &
    Invalid or rejected actions show clear system responses, and no improper
    content modification path is exposed through normal user actions. &
    Screenshots, error output &
    Passed / Partially Completed / Failed \\
    \hline

    UAT-06 &
    TC-REG-003, TC-REG-005 &
    Manage related student or integrator records. &
    Student or integrator information can be viewed and managed correctly, and
    related records remain visible and consistent in the interface. &
    Screenshots, updated records &
    Passed / Partially Completed / Failed \\
    \hline

  \end{tabular}
  }%
  \captionsetup{type=table}
  \caption{User acceptance success criteria for Tenant Administrator}
  \label{tab:uat-success-tenant}
\end{table}

% =========================================================
% IT Integrator UAT
% =========================================================
\begin{table}[H]
  \centering
  \smalltable
  \scriptsize
  {%
  \setlength{\tabcolsep}{3pt}
  \renewcommand{\arraystretch}{1.18}
  \begin{tabular}{|>{\raggedright\arraybackslash}p{0.09\textwidth}|
                  >{\raggedright\arraybackslash}p{0.15\textwidth}|
                  >{\raggedright\arraybackslash}p{0.22\textwidth}|
                  >{\raggedright\arraybackslash}p{0.30\textwidth}|
                  >{\raggedright\arraybackslash}p{0.14\textwidth}|
                  >{\centering\arraybackslash}p{0.10\textwidth}|}
    \hline
    \textbf{UAT ID} &
    \textbf{Linked TC ID(s)} &
    \textbf{Task} &
    \textbf{Success criteria for passed result} &
    \textbf{Evidence} &
    \textbf{Result} \\
    \hline\hline

    UAT-07 &
    TC-ITO-001, TC-ITO-002, TC-ITO-003, TC-ITO-004, TC-ITO-005 &
    Execute the intended API integration workflow. &
    Valid API requests return the expected identifiers, metadata, download, or
    status results, while invalid requests are rejected without creating an
    unintended record. &
    API responses, logs &
    Passed / Partially Completed / Failed \\
    \hline

  \end{tabular}
  }%
  \captionsetup{type=table}
  \caption{User acceptance success criteria for IT Integrator / Developer}
  \label{tab:uat-success-integrator}
\end{table}

% =========================================================
% External Verifier UAT
% =========================================================
\begin{table}[H]
  \centering
  \smalltable
  \scriptsize
  {%
  \setlength{\tabcolsep}{3pt}
  \renewcommand{\arraystretch}{1.18}
  \begin{tabular}{|>{\raggedright\arraybackslash}p{0.09\textwidth}|
                  >{\raggedright\arraybackslash}p{0.15\textwidth}|
                  >{\raggedright\arraybackslash}p{0.22\textwidth}|
                  >{\raggedright\arraybackslash}p{0.30\textwidth}|
                  >{\raggedright\arraybackslash}p{0.14\textwidth}|
                  >{\centering\arraybackslash}p{0.10\textwidth}|}
    \hline
    \textbf{UAT ID} &
    \textbf{Linked TC ID(s)} &
    \textbf{Task} &
    \textbf{Success criteria for passed result} &
    \textbf{Evidence} &
    \textbf{Result} \\
    \hline\hline

    UAT-08 &
    TC-STU-001, TC-STU-003, TC-EXT-001, TC-EXT-002, TC-E2E-002, TC-E2E-003 &
    Verify a report through the public portal using QR code or CID entry. &
    The public portal displays the correct result for valid, invalid,
    tampered, or unavailable cases without requiring internal access and
    without exposing sensitive internal details. &
    Screenshots, scan proof, portal output &
    Passed / Partially Completed / Failed \\
    \hline

  \end{tabular}
  }%
  \captionsetup{type=table}
  \caption{User acceptance success criteria for External Verifier}
  \label{tab:uat-success-verifier}
\end{table}

During user testing, each participant will receive a task script derived from
the UAT table corresponding to the participant's assigned role. The script
will describe the initial state, the assigned task, and the expected result.
Observers will record task completion, time taken, defects encountered, and
any task blockers observed during execution.

User acceptance testing will be considered satisfactory when assigned tasks
are completed successfully, expected results are observed, critical defects
are resolved, and regression re-testing confirms that fixes do not introduce
new failures. At least 90\% of critical UAT tasks should be completed
without assistance for each user group before acceptance is concluded.

\subsubsection{User acceptance result classification}

For user acceptance testing, each scenario will be recorded using one of the
following outcome labels.

\begin{table}[H]
  \centering
  \smalltable
  \footnotesize
  {%
  \setlength{\tabcolsep}{6pt}
  \renewcommand{\arraystretch}{1.3}
  \begin{tabular}{|>{\raggedright\arraybackslash}p{0.24\textwidth}|
                  >{\raggedright\arraybackslash}p{0.66\textwidth}|}
    \hline
    \textbf{Outcome label} & \textbf{Definition} \\
    \hline\hline

    \textbf{Passed} &
    The participant completes the assigned task successfully and the expected
    result is observed. \\
    \hline

    \textbf{Partially Completed} &
    The participant completes part of the assigned task, but encounters minor
    defects, incomplete behaviour, or non-blocking issues during execution. \\
    \hline

    \textbf{Failed} &
    The participant cannot complete the assigned task, or the observed result
    does not satisfy the expected result for the intended role. \\
    \hline

  \end{tabular}
  }%
  \captionsetup{type=table}
  \caption{User acceptance test result classification}
  \label{tab:uat-result-classification}
\end{table}

These outcome labels will be recorded together with task
evidence, and any defects identified during the session.

\subsection{Security Testing}

Security evaluation for VERACISM will focus on protecting the
confidentiality, integrity, and availability of data in line with ISM
security policies. In particular, the system is expected to meet the
following security targets:

\begin{itemize}
  \item All network traffic uses HTTPS with current TLS versions;
  \item Passwords and secrets are never logged or exposed in client responses;
  \item Input validation and output encoding are applied to all external interfaces;
  \item The system is regularly scanned for vulnerabilities and no critical or high
        severity findings remain unresolved beyond agreed timelines.
\end{itemize}

VERACISM will be subjected to Vulnerability Assessment and Penetration Testing
(VAPT) before the initial go-live and after any significant change that may
affect its security posture. The target of the assessment is the VERACISM web
application and APIs, reachable under the staging URL
\texttt{https://sb-[update-domain-if-needed]} (to be confirmed). Test
credentials will be provisioned securely (for example, via encrypted email or
a password manager). Systems, subdomains, and third-party services not
explicitly included in the scope will remain out of scope.

\begin{table}[htbp]
  \centering
  \smalltable
  \footnotesize
  {%
  \setlength{\tabcolsep}{8pt}
  \renewcommand{\arraystretch}{1.3}
  \begin{tabular}{|>{\raggedright\arraybackslash}p{3.6cm}|
                  >{\raggedright\arraybackslash}p{4.4cm}|
                  >{\raggedright\arraybackslash}p{4.4cm}|}
    \hline
    \textbf{Planned security activity} &
    \textbf{Tools (examples)} &
    \textbf{Objective / scope} \\
    \hline\hline

    Automated vulnerability scanning &
    Burpsuite &
    Discover CVEs, open ports and services, misconfigurations, weak defaults,
    outdated components, and common web issues such as XSS, SQL injection,
    CSRF, and insecure headers. \\
    \hline

    Dynamic application scan (DAST) &
    Burpsuite &
    Scan the VERACISM web application and APIs (login, upload, verification
    portal, internal console) under authenticated and unauthenticated sessions. \\
    \hline

    Manual verification of findings &
    Manual review using browser tools and scripts &
    Remove false positives from automated scans and assess real impact and
    exploitability of each issue in the context of VERACISM. \\
    \hline

    Manual penetration testing &
    Custom payloads and test harnesses aligned with OWASP Top 10 &
    Perform targeted testing of high-risk flows and endpoints, based on the
    focus areas in Table~\ref{table:owasp-focus}. \\
    \hline

    Continuous dependency and image scanning (post-deployment) &
    Snyk or an ISM-approved equivalent integrated in CI/CD &
    Detect vulnerable libraries, container images, and IaC templates over time
    and feed findings into ISM's vulnerability management process. \\
    \hline
  \end{tabular}
  }%
  \captionsetup{type=table}
  \caption{Planned security testing activities for VERACISM}
  \label{table:security-testing-plan}
\end{table}

\begin{table}[htbp]
  \centering
  \smalltable
  \footnotesize
  {%
  \setlength{\tabcolsep}{6pt}
  \renewcommand{\arraystretch}{1.3}
  \begin{tabular}{|>{\raggedright\arraybackslash}p{1.5cm}|
                  >{\raggedright\arraybackslash}p{4.0cm}|
                  >{\raggedright\arraybackslash}p{6.9cm}|}
    \hline
    \textbf{OWASP ID} &
    \textbf{Risk name} &
    \textbf{Planned focus in VERACISM} \\
    \hline\hline

    A01 &
    Broken Access Control &
    Verify RBAC/ABAC enforcement, least privilege for service accounts, and
    strict access restrictions on reports, admin pages, and APIs. \\
    \hline

    A02 &
    Cryptographic Failures &
    Check HTTPS/TLS configuration, key and secret handling, and correct usage
    of hashing (SHA-256) and encryption where required. \\
    \hline

    A03 &
    Injection &
    Confirm use of parameterized queries and ORM; test all inputs for SQL,
    command, and template injection. \\
    \hline

    A04 &
    Insecure Design &
    Review threat model, design of immutability and auditability, and presence
    of required security controls in high-risk flows. \\
    \hline

    A05 &
    Security Misconfiguration &
    Test for insecure defaults, missing security headers, verbose error messages,
    open ports, and weak configuration in .NET, MSSQL, and hosting. \\
    \hline

    A06 &
    Vulnerable and Outdated Components &
    Confirm that dependency scans are running and that findings in .NET,
    Angular/React, and IPFS/Filecoin libraries are addressed. \\
    \hline

    A07 &
    Identification and Authentication Failures &
    Test login flows, token handling, session timeout, brute-force resistance,
    and integration with the ISM identity provider. \\
    \hline

    A08 &
    Software and Data Integrity Failures &
    Validate CI/CD controls, integrity of deployed artifacts, and that report
    hashes and CIDs reliably detect unauthorized changes. \\
    \hline

    A09 &
    Security Logging and Monitoring Failures &
    Confirm that authentication, access control, upload, download, and
    verification events are logged and integrated with monitoring/alerting. \\
    \hline

    A10 &
    Server-Side Request Forgery (SSRF) &
    Test outbound calls (Lighthouse, IPFS gateways, explorers) against
    allow-lists and ensure no user-controlled URLs are used. \\
    \hline
  \end{tabular}
  }%
  \captionsetup{type=table}
  \caption{OWASP Top 10 focus areas for VERACISM penetration testing}
  \label{table:owasp-focus}
\end{table}

\clearpage
The overall VAPT procedure for VERACISM will follow these steps:

\begin{enumerate}
  \item \textbf{Information gathering:}  
        Gather information on infrastructure, technology stack, exposed services, and endpoints; Establish an attack surface for web application and APIs.

  \item \textbf{Automated vulnerability scanning:}  
           Perform network and application scans by using the tools in
        Table~\ref{table:security-testing-plan}. The aim is to recognize:
        \begin{itemize}
          \item known systems and libraries with recent CVEs;
          \item open ports and exposed services;
          \item misconfigurations and weak defaults;
          \item outdated software components;
          \item exposure in the web due to XSS, SQL injection, CSRF, missing or insecure HTTP headers.
        \end{itemize}

  \item \textbf{Manual verification of findings:}  
        Check the results of the automation process manually, in order to eliminate false positives and to gain insight into the functionality of each error in a running VERACISM instance.

  \item \textbf{Manual security testing:}  
        Conduct regionally specific testing of affected flows and endpoints with safe, non-destructive payloads. Take screenshots and logs of any evidence for each confirmed find, record severity, and note reproduction steps, and identify payloads.

  \item \textbf{Penetration testing focus areas:}  
        A Use the OWASP Top 10 mean areas outlined in
        Table~\ref{table:owasp-focus} to test each of the areas, including access control, crypto, input handling, configuration, components, logging, and outbound calls, thoroughly.

  \item \textbf{Remediation and re-testing:}  
        Deliver comprehensive, technical conclusions and remediation tips for the development and ops teams. Re-test following fixes to ensure vulnerabilities addressed and remediation did not create new vulnerabilities.

  \item \textbf{Reporting:}  
        Generate a formal VAPT report that outlines all of the verified vulnerabilities, their severity and potential impact, reproduction and suggested remediation actions. Document and retain reports and re-test results for audits and reevaluation.
\end{enumerate}

Optionally, security assessment can be conducted along with load testing performed using Azure Load Testing or other load testing tools. In this mode, realistic user load to be applied to VERACISM while observing security controls, rate limiting, throttling and brute forced protection under stress.

Continuous monitoring of VERACISM is recommended by using dependency and image scanning tools to help maintain the security of VERACISM over time. The results of these multi-sensors will be managed as part of ISM's vulnerability management process.

Security acceptance criteria for the project are:

\begin{itemize}
  \item All Critical and High vulnerabilities identified in VAPT or automated
        scans must be fixed or formally risk-accepted by ISM Information
        Security before production release;
  \item Medium and Low issues must have a documented remediation plan and
        target date;
  \item VAPT reports, scan results, and re-test outcomes must be retained as
        evidence for audits and internal reviews.
\end{itemize}
